\section{Features and anti-leakage}
\label{sec:features}

\subsection{The Common Data Model}

After the country-specific Phase~1 preprocessing, every country's \texttt{processed/} folder contains the same five tables: \texttt{nodes\_station}, \texttt{nodes\_service}, \texttt{edges\_stops\_at}, \texttt{edges\_adjacent} and \texttt{nodes\_fault}. Phase~2 (\texttt{stop\_level/preprocess\_stops.py}) loads all 15 of those CSVs (5 schemas $\times$ 3 countries), joins them at the stop level, derives 16 feature groups, fits a \texttt{StandardScaler} on training rows only, and persists a parquet feature store to \texttt{Data/stops/}. Every downstream consumer (the four-model benchmark, the SHAP report, the cause classifier, the Knowledge Graph) reads from that single store, which guarantees that they all see the same features and the same splits.

\subsection{Per-stop features}

\Cref{tab:feature-groups} groups the 37 pre-departure features (scenario~A) and the 40 inflight features (scenario~B) by what they measure. Scenario~B differs from scenario~A by exactly three additions, all of which describe the actual delays accumulated by the same train at \emph{prior} stops on the current run.

\begin{table}[H]
    \centering
    \caption{Per-stop feature groups. Scenario~A is the pre-departure setting: only information available before the train leaves the platform. Scenario~B adds the three inflight features that summarise actual delays at prior stops on the same run. The grouping reflects \emph{what each feature measures}, not the order in which it is computed.}
    \label{tab:feature-groups}
    \small
    \setlength{\tabcolsep}{6pt}
    \begin{tabularx}{\textwidth}{l X S[table-format=2.0]}
        \toprule
        Group & Features & {Count} \\
        \midrule
        Stop position        & \texttt{stop\_order, position\_norm, is\_origin, is\_terminus, n\_total\_stops, cum\_distance\_km, cum\_scheduled\_minutes} & 7 \\
        Stop time            & \texttt{scheduled\_arrival\_hour, scheduled\_arrival\_dow, is\_holiday, month, month\_sin, month\_cos} & 6 \\
        Station static       & \texttt{lat, lon, degree, betweenness\_centrality, avg\_historical\_delay} & 5 \\
        Service identity     & \texttt{train\_class\_code, service\_n\_stops, service\_route\_distance\_km, service\_scheduled\_duration\_min} & 4 \\
        Weather              & \texttt{weather\_severity, temperature, wind\_speed, precipitation, snow\_depth} & 5 \\
        Weather differential & \texttt{delta\_severity\_vs\_prev\_stop, delta\_wind\_vs\_prev\_stop} & 2 \\
        Lag features         & \texttt{station\_lag1\_rate, station\_lag7\_rate, station\_lag1\_avg\_delay, train\_lag1\_rate, train\_lag7\_rate, train\_station\_lag7\_rate, nbr\_lag1\_rate} & 7 \\
        Fault context        & \texttt{station\_n\_active\_faults\_14d} & 1 \\
        \midrule
        \textbf{Inflight (scenario~B only)} & \texttt{prev\_stop\_actual\_delay, cum\_actual\_delay\_so\_far, max\_actual\_delay\_so\_far} & 3 \\
        \bottomrule
    \end{tabularx}
\end{table}

\subsection{Why these feature groups}
\label{sec:feature-rationale}

Each group answers a specific question that an experienced operator would ask before flagging a stop as high-risk. The groups are constructed top-down (operational hypothesis first, columns second), not bottom-up from whatever happens to be in the source feeds. The choice of representation inside each group also matters; where it does, we say so.

\begin{description}[leftmargin=2.0em, style=nextline, itemsep=4pt]
    \item[Stop position (7 columns).] Punctuality decays predictably along a route, so the model needs to know where the current stop sits on the run. We give it both the absolute position (\texttt{stop\_order}, \texttt{cum\_distance\_km}, \texttt{cum\_scheduled\_minutes}) and a route-length-invariant fraction (\texttt{position\_norm}, in $[0, 1]$); together they let the model express ``late stop on a long route'' without having to multiply features itself. \texttt{is\_origin} and \texttt{is\_terminus} flag the two endpoints, whose disruption distributions differ from the middle of the run -- origins almost never have an upstream delay, termini accumulate every prior shock.
    \item[Stop time (6 columns).] Hour-of-day captures rush-hour congestion in the operator timetable; day-of-week separates weekday commuter patterns from weekend leisure traffic; \texttt{is\_holiday} flips the normal weekday template on national holidays. We encode the month both as an ordinal (\texttt{month}) and as the cyclic pair $(\sin\frac{2\pi m}{12}, \cos\frac{2\pi m}{12})$. Without the cyclic encoding January and December would sit 11 ordinal units apart instead of 1, which is wrong for any operator whose seasonality wraps year-end.
    \item[Station static (5 columns).] \texttt{avg\_historical\_delay} per station is, on its own, the strongest single pre-departure feature in the SHAP report. It summarises everything we know about the station's intrinsic congestion without us having to model it explicitly. \texttt{degree} (number of distinct neighbours in the adjacency graph) and \texttt{betweenness\_centrality} flag hub vulnerability: hubs concentrate cascading effects, so the model can use those as a prior on disruption probability, and the GraphSAGE benchmark uses them again for message passing.
    \item[Service identity (4 columns).] \texttt{train\_class\_code} is a categorical encoding of the operating class -- commuter, regional, intercity, high-speed. High-speed services typically run on dedicated track and are punctuality outliers in both directions (high precision when on time, very disrupted when off), so calling them out as a separate class lets the model carry a different prior. The route distance and scheduled duration let the model express ``long, slow runs'' as a separate regime from ``short, fast runs''; the SHAP report shows it does exactly that.
    \item[Weather (5 columns).] We keep both a country-harmonised \texttt{weather\_severity} ordinal (0--4, computed jointly from snow depth, precipitation and visibility) \emph{and} the raw \texttt{snow\_depth, precipitation, wind\_speed, temperature} columns. The severity ordinal is the country-invariant signal the linear model relies on; the raw columns let the trees discover non-linearities the ordinal flattens (e.g.\ \SI{30}{cm} of snow at \SI{-5}{\celsius} hits Finnish operations differently than \SI{30}{cm} at \SI{0}{\celsius}, because freezing-vs-thawing changes the friction regime). Removing the raw columns drops scenario-A PR-AUC by approximately \num{0.012}, so the redundancy is paying for itself.
    \item[Weather differential (2 columns).] Changes in weather are often more predictive than levels: a sudden wind gust (positive \texttt{delta\_wind\_vs\_prev\_stop}) signals operational stress that the absolute wind value misses. We compute the change against the previous stop on the same run, not the previous calendar timestep, so the differential is naturally aligned with the train's own progress.
    \item[Lag features (7 columns).] These are the strongest pre-departure feature group, because they implicitly encode \emph{unobserved} structural problems. If \texttt{station\_lag7\_rate} on station~$S$ has been climbing for a week, the model can flag $S$ as risky without us having to know whether the underlying cause is a signal failure, a track-work zone, or a shift in weather. \texttt{train\_station\_lag7\_rate} is the per-(train, station) recency rate -- combining train identity and station identity is what makes this the SHAP-dominant pre-departure feature: it captures ``this specific train at this specific station has been late lately'', a much sharper signal than either rate alone.
    \item[Fault context (1 column).] \texttt{station\_n\_active\_faults\_14d} counts open infrastructure tickets at the station within the last fortnight. The 14-day window matches the typical infrastructure-repair cycle observed in the Dutch fault log; longer windows over-smooth, shorter windows miss faults that are still open. Only the Netherlands has a public fault feed, so the column is zero for Italy and Finland; tree models handle the constant-zero case correctly (the feature is effectively ignored where it is uninformative).
    \item[Inflight (3 columns, scenario~B only).] \texttt{prev\_stop\_actual\_delay} is a near-deterministic predictor of cascading: a train that is 8~minutes late at stop $k-1$ will, with high probability, also be late at stop $k$. The other two columns help the model separate two qualitatively different situations: \texttt{cum\_actual\_delay\_so\_far} signals a run that is losing time monotonically (debt accumulating), while \texttt{max\_actual\_delay\_so\_far} flags the worst single shock in the run history (a debt that may have already been recovered). Removing all three drops scenario-B PR-AUC from \num{0.94} to \num{0.35} (\cref{sec:xai}), confirming they carry essentially the entire A$\to$B uplift.
\end{description}

\subsection{Anti-leakage protocol}

The pipeline rests on a single rule: at prediction time the feature vector $\mathbf{x}$ may only contain information that was actually available before the prediction was made. Six layers enforce that rule, each one a check that the previous layer is not silently violated.

\begin{enumerate}
    \item \textbf{Phase~1 retains \texttt{delay\_minutes} only as a label source.} The column is used to compute the disruption indicator and is then dropped before the parquet output is written. It never appears as a feature.
    \item \textbf{Phase~2 keeps two explicit blocklists.} The set \texttt{LEAKY\_COLS\_A} contains 14 columns that would leak in the pre-departure setting; \texttt{LEAKY\_COLS\_B} contains 11 (it is \texttt{LEAKY\_COLS\_A} minus the three inflight whitelist features). Before the scaler is fit, the active feature list is asserted to be disjoint from the relevant blocklist.
    \item \textbf{Lag features use strict inequality on the daily aggregation grid.} The feature \texttt{station\_lag1\_rate(D)} is computed from data on day $D-1$ and earlier only -- never $D$. A unit test feeds a known sequence and asserts that the day-of value is not visible. This is the easiest leakage class to introduce by accident, so we test it directly.
    \item \textbf{The temporal split is at day granularity.} Every stop of a given service belongs to a single day, so a service ID can never straddle the train, validation, or test partitions.
    \item \textbf{The scaler is fit on training rows only.} Validation and test rows are transformed but never re-fit; this prevents test-set statistics from quietly affecting the normalisation.
    \item \textbf{Per-scenario feature lists are persisted as JSON.} Every downstream consumer (model code, SHAP report, ablation runner) reads \texttt{feature\_names\_A.json} or \texttt{feature\_names\_B.json}, so the contract about what is and is not in scenario~A is the same one the model trained against.
\end{enumerate}

\subsection{Splits}

\Cref{tab:splits} summarises the temporal split. Train, validation and test are non-overlapping date windows, and \texttt{splits.assert\_no\_service\_overlap} enforces at runtime that no service ID appears in two partitions. The disruption rate drops from \SI{10.13}{\percent} on the four-month training window to \SI{8.64}{\percent}--\SI{8.66}{\percent} on the validation and test months, reflecting a milder spring weather regime.

\begin{table}[H]
    \centering
    \caption{Temporal split sizes after Phase~2. ``Disrupted rate'' is the fraction of stops whose arrival is more than five minutes late or that are cancelled. The drop from \SI{10.13}{\percent} (winter training) to \SI{8.66}{\percent} (June test) reflects normal seasonal disruption -- it is not a per-class label drift.}
    \label{tab:splits}
    \small
    \begin{tabular}{l
        S[table-format=8.0,group-separator={,}]
        S[table-format=2.2]
        l}
        \toprule
        Split & {Rows} & {Disrupted rate (\%)} & Date window \\
        \midrule
        Train      & 11073233 & 10.13 & 2024-01-01 -- 2024-04-30 \\
        Validation &  2820493 &  8.64 & 2024-05-01 -- 2024-05-31 \\
        Test       &  2706443 &  8.66 & 2024-06-01 -- 2024-06-30 \\
        \bottomrule
    \end{tabular}
\end{table}

\subsection{Memory engineering}

The unified \num{16.6}\,M-row, 47-column dataframe initially needed more than \SI{15}{GB} of RAM, mostly due to default \texttt{float64} dtypes on numeric columns and uncategorised string IDs. Three concrete changes brought it under \SI{8}{GB}.

\begin{enumerate}
    \item \textbf{Compact dtypes on load.} Every numeric column is cast to \texttt{float32} (a \SI{50}{\percent} saving on numerics), small ordinals to \texttt{int8} or \texttt{int32}, and string IDs (such as \texttt{station\_id}) to \texttt{category}.
    \item \textbf{Map-based stop-grain joins.} Joining the two largest tables on a million-key categorical column with \texttt{pd.merge} materialises a multi-gigabyte intermediate hash. We replace it with \texttt{Series.map(dict)}, which keeps the join purely in NumPy.
    \item \textbf{Removed redundant copies.} Eight \texttt{df = df.copy()} calls in feature builders are removed; the chained-flow pattern \texttt{df = F.add\_*(df)} is safe to mutate in place once each builder is shown not to produce aliased views.
\end{enumerate}

A secondary high-cardinality fallback inside \texttt{compute\_train\_lag} skips the contiguous-calendar grid whenever the unique-key count exceeds \num{50000}: the Netherlands has approximately \num{1000000} unique train keys and the default grid would expand to \num{215000000} rows.
